\chapter{数字经济与经济增长}

第 7 章末尾提出的两个问题构成本章的任务：以增长核算为标尺，检验数据要素是否真的改变了增长的路径，并解开数字技术留下的生产率谜题。本章先回顾增长核算方法与索洛稳态（8.1 节），再正面处理索洛悖论（8.2 节），随后把数据要素作为独立生产要素放进增长模型（8.3 节），最后讨论数字基础设施的外部性与增长收敛（8.4 节）。

\begin{examguide}
\textbf{本章考情}：增长核算与索洛模型（含黄金律水平）是学硕宏观必考的经典内容（\important{学硕·热点}），数字经济背景下的"索洛悖论"与数据要素增长模型是 0258 论述题高频素材（\important{专硕·核心}）。增长核算方程的推导与贡献率分解是计算题的固定命题形式。

\textbf{核心考点}：（1）增长核算方程的对数微分推导与索洛余值；（2）欧拉定理、收入份额与增长分解数值算例；（3）索洛稳态与黄金律水平；（4）索洛悖论与三种解释；（5）消费者剩余测度之辨与 GDP-B；（6）鲍莫尔病与数字化的三条出路；（7）数据要素三要素核算与规模报酬递增；（8）AK 模型与数据要素的内生增长含义；（9）数据产权配置与数据积累稳态；（10）数字基础设施的外部性与条件收敛。

\textbf{本章主线}：如何核算增长（8.1）→ 数字技术为何"看不见"（8.2）→ 数据如何改变增长（8.3）→ 基础设施与收敛（8.4）。
\end{examguide}

\section{增长核算方法}

\subsection{增长核算方程}

增长核算（growth accounting）将产出增长分解为要素投入增长与全要素生产率（TFP）增长，是数字经济宏观分析的基准工具。\mn{增长核算方程的用途（考试高频）：（1）计算 TFP 增长率（索洛余值）；（2）分解各要素对增长的贡献率（$\alpha \dot K/K$ 除以 $\dot Y/Y$）；（3）判断增长方式（要素驱动 vs 效率驱动），中国语境下的标准表述是从"要素驱动"转向"创新驱动"的增长方式转变。}无论是理解"创新驱动"增长方式的度量基础，还是检验数字技术是否兑现增长承诺（8.2 节），这一步分解都是分析的起点。

\begin{definition}[全要素生产率（TFP）]
全要素生产率是产出水平与加权要素投入水平之比 $A = Y/(K^\alpha L^{1-\alpha})$，度量给定要素投入下的产出水平高低，即各要素综合使用效率的水平；技术进步、组织改进与资源配置改善都沉淀在其中。注意区分：TFP 是\textbf{水平}概念，增长核算反解出的"余值"是它的\textbf{增长率}，即下文定义的索洛余值。
\end{definition}

把这个水平概念的变化率单独分离出来，就得到增长核算的核心结果。衡量效率变化快慢的指标，正是下面定义的索洛余值。

\begin{definition}[索洛余值（Solow residual）]
索洛余值指产出增长率中无法由要素投入增长率解释的部分，即 TFP 的增长率
\[
\frac{\dot A}{A} = \frac{\dot Y}{Y} - \alpha\frac{\dot K}{K} - (1-\alpha)\frac{\dot L}{L},
\]
因需要从增长核算方程中"反解"（作差扣除）得到而得名。
\end{definition}

增长核算要回答的是\textbf{归因}问题：今年的产出增长中，多少来自多投了资本、多少来自多雇了工人、多少来自效率提升？生产函数 $Y = AK^\alpha L^{1-\alpha}$ 写的是水平之间的关系，无法直接读出增长率层面的答案，必须把它改写为增长率形式。对数微分正是完成这一改写的手段：$\ln(uv)=\ln u+\ln v$ 把乘积关系化为加和关系，$\frac{d\ln X}{dt}=\frac{\dot X}{X}$ 把求导直接写成增长率，两者结合，产出增长才能被干净地拆成各来源之和。

\begin{derivation}{增长核算方程}{}
\textbf{增长核算方程的推导。}设总量生产函数为科布—道格拉斯形式：
\[
Y = A K^\alpha L^{1-\alpha},
\]
其中 $A$ 为全要素生产率，$\alpha$ 与 $1-\alpha$ 分别为资本与劳动的产出弹性（在竞争性要素市场假设下等于各自的收入份额）。第一步，两边取对数，把乘积关系化为加和关系：
\[
\ln Y = \ln A + \alpha\ln K + (1-\alpha)\ln L.
\]
第二步，对时间求导（对数微分法），利用 $\frac{d\ln X}{dt} = \frac{1}{X}\frac{dX}{dt}$ 把每个对数项写成自身的增长率：
\[
\frac{d\ln Y}{dt} = \frac{d\ln A}{dt} + \alpha\frac{d\ln K}{dt} + (1-\alpha)\frac{d\ln L}{dt}.
\]
第三步，记增长率 $\dot X/X = d\ln X/dt$，得\important{增长核算方程}：
\[
\frac{\dot Y}{Y} = \frac{\dot A}{A} + \alpha\frac{\dot K}{K} + (1-\alpha)\frac{\dot L}{L}.
\]
即
\[
\text{产出增长率} = \text{TFP 增长率} + \alpha\cdot\text{资本增长率} + (1-\alpha)\cdot\text{劳动增长率}.
\]
\end{derivation}

余值度量的是要素使用效率的综合提升：技术进步、组织改进、规模经济、资源配置改善，一切不体现为要素数量增加的增长来源都沉淀在 $A$ 中。测算中产出弹性 $\alpha$ 的取值通常直接使用要素收入份额：竞争性要素市场下，要素按边际产品取得报酬，资本报酬率 $r=\partial Y/\partial K=\alpha Y/K$，故资本收入份额 $rK/Y=\alpha$，劳动同理。这一对应关系由\textbf{欧拉定理}（Euler's theorem）保证：生产函数满足规模报酬不变（constant returns to scale，$\alpha+(1-\alpha)=1$）时，各要素按边际产品取得的报酬之和恰好等于总产出，$(\partial Y/\partial K)K + (\partial Y/\partial L)L = \alpha Y + (1-\alpha)Y = Y$，产出弹性因此可以直接用收入份额校准。

以下示例数据演示分解计算：设某年 $\dot Y/Y = 6\%$、$\dot K/K = 5\%$、$\dot L/L = 0.5\%$，取资本收入份额 $\alpha = 0.4$，则索洛余值为 $\dot A/A = 6\% - 0.4\times 5\% - 0.6\times 0.5\% = 3.7\%$；三类来源对增长的贡献率分别为：资本 $0.4\times 5\%/6\% \approx 33.3\%$、劳动 $0.6\times 0.5\%/6\% = 5\%$、TFP $3.7\%/6\% \approx 61.7\%$，在示例设定下增长的主要来源是效率提升。现实测算中，2023 年数字经济全要素生产率对中国经济增长的整体贡献约为 22.5\%，在全要素生产率对经济增长的贡献中占比过半（中国信息通信研究院，2024）。

\important{增长核算的意义在于把"增长"这一总量概念分解为可归因的成分}；其局限同样明显：余值是"残差"，其中混杂了测度误差与模型设定误差，数字经济时代这一局限被显著放大（8.2 节）。

\subsection{索洛模型的稳态（回顾）}

增长核算回答"增长由什么实现"，而判断增长的长期前景需要回答"增长由什么驱动"，后者正是增长理论的核心问题。增长理论要解释的是一国人均收入为什么长期增长、国与国之间的收入差距为什么持续存在。本章涉及的模型构成一条完整的回答链条：索洛模型把长期增长归结为外生技术进步（本节回顾稳态与黄金律水平）；8.3 节讨论数据要素能否把这个"外生"内生化；8.4 节的收敛方程也直接由本节导出。索洛模型的关键结论可以先行概括：人均产出的长期增长完全来自技术进步率 $g$，要素积累决定收入水平而不决定增长速度，后续各节都建立在这一结论之上。

\begin{derivation}{索洛稳态}{}
\textbf{稳态推导。}设储蓄率为 $s$，资本折旧率为 $\delta$，人口增长率为 $n$，生产函数为 $Y = K^\alpha (AL)^{1-\alpha}$（劳动增进型技术进步，labor-augmenting technological progress，技术进步率 $g$）。人均有效资本 $k = K/(AL)$ 的动态方程为
\[
\dot k = s k^\alpha - (\delta + n + g) k.
\]
稳态要求 $\dot k = 0$，即 $s (k^*)^\alpha = (\delta+n+g) k^*$，整理解得
\[
k^* = \left(\frac{s}{\delta + n + g}\right)^{\frac{1}{1-\alpha}},
\]
稳态人均有效产出 $y^* = (k^*)^\alpha$。
\end{derivation}

图~\ref{fig:solowphase} 给出了稳态的直观理解：储蓄是资本存量的"入水"，折旧、人口增长与技术进步构成三个"出水口"；入水线 $sk^\alpha$ 随 $k$ 增加而弯曲放缓（资本边际收益递减），出水线 $(\delta+n+g)k$ 随 $k$ 线性上升，两者相等处就是稳态 $k^*$，任何偏离都会被拉回（$k<k^*$ 时 $\dot k>0$，$k>k^*$ 时 $\dot k<0$）。稳态的含义：稳态下 $k$ 与 $y$ 都不再增长，人均产出 $Y/L = A y^*$ 的增长完全来自 $A$ 的增长，因此\important{人均产出的长期增长率完全由技术进步率 $g$ 决定}：没有技术进步，人均增长最终归于停滞。\mn{稳态公式速记：$k^*=\left(\frac{s}{\delta+n+g}\right)^{\frac{1}{1-\alpha}}$，上储蓄、下三耗、指数看资本份额；长期增长率只认 $g$ 不认 $s$。}这一结论为理解数字经济的增长角色提供了基准：\important{数字技术的增长贡献，最终体现为它对 $g$（技术进步率）与 $A$（效率水平）的影响}；数字技术若只是改变了资本形态而未提升 $g$，则其长期增长效应有限。附带一提两种写法的对应：把本节生产函数改写成 8.1.1 的核算形式 $Y = \tilde A K^\alpha L^{1-\alpha}$ 时，$\tilde A = A^{1-\alpha}$，稳态下核算余值的增长率为 $(1-\alpha)g$ 而非 $g$；差别只在于技术进步记在了哪里，不影响任何经济结论。

\begin{figure}[htbp]
\centering
\begin{tikzpicture}[>=Stealth, scale=1]
\draw[->] (0,0) -- (7.2,0) node[below left=8pt] {$k$};
\draw[->] (0,0) -- (0,3.3);
\draw[second,thick] (0,0) -- (6.2,2.319);
\draw[main,thick] plot[domain=0:6.2,samples=100] (\x,{0.7*pow(\x,0.5)});
\draw[dashed,structurecolor] (3.5,0) -- (3.5,1.31);
\fill (3.5,1.31) circle [radius=1.6pt];
\node[below] at (3.5,0) {$k^*$};
\node[font=\small] at (5.45,1.14) {流入 $sk^\alpha$};
\node[font=\small] at (4.7,2.56) {流出 $(\delta+n+g)k$};
\draw[->,thick] (0.5,0.06) -- (2.4,0.06);
\draw[->,thick] (5.9,0.06) -- (4.3,0.06);
\node[font=\small] at (1.45,-0.35) {$\dot k>0$};
\node[font=\small] at (5.1,-0.35) {$\dot k<0$};
\end{tikzpicture}
\caption{索洛模型的稳态}
\label{fig:solowphase}
\end{figure}

稳态回答了长期增长由什么驱动，但没有回答稳态本身是否划算：同样的技术进步率下，不同的储蓄率对应不同的稳态消费水平。使稳态人均消费最大化的资本水平称为\textbf{黄金律水平}（golden rule），它同时给出储蓄率的效率边界。

\begin{derivation}{黄金律水平}{}
\textbf{黄金律水平的推导。}稳态人均消费等于产出减去投资：
\[
c^* = f(k^*) - (\delta+n+g)k^*.
\]
稳态资本 $k^*$ 随储蓄率 $s$ 单调上升，最大化稳态消费等价于对 $k^*$ 求一阶条件，令导数为零：
\[
f'(k_{GR}) = \delta+n+g.
\]
在科布—道格拉斯形式 $f(k) = k^\alpha$ 下，$f'(k) = \alpha k^{\alpha-1}$，解得
\[
k_{GR} = \left(\frac{\alpha}{\delta+n+g}\right)^{\frac{1}{1-\alpha}}.
\]
与稳态水平 $k^* = (s/(\delta+n+g))^{1/(1-\alpha)}$ 比较：$k^* > k_{GR}$ 当且仅当 $s > \alpha$。
\end{derivation}

黄金律条件 $f'(k_{GR})=\delta+n+g$ 的经济含义：资本边际产品恰好补偿三种损耗（折旧、人口稀释与技术进步稀释），再多积累一单位资本的净收益为零。$s<\alpha$ 时稳态资本低于黄金律水平，储蓄不足，提高储蓄率可以提高稳态消费，但要以牺牲当期消费为代价；$s>\alpha$ 时资本过度积累，出现\textbf{动态无效率}（dynamic inefficiency）：降低储蓄率能同时提高当期消费与稳态消费，所有世代的境况都可以变好。黄金律水平与动态无效率的判断是学硕宏观的常考概念（辨析题常客）：储蓄率并非越高越好。

\section{索洛悖论}

\subsection{悖论的提出}

索洛悖论诞生于 1980 年代美国"生产率放缓"的学术争论。1970 年代起，美国劳动生产率增速从战后二十余年的年均近 3\% 滑落至 1.5\% 以下，围绕放缓成因的争论持续了十余年；与此同时，计算机等信息技术设备的投资高速扩张，本应成为生产率新引擎的技术革命，在统计上却迟迟不见踪影。1987 年，索洛在《纽约时报》书评中留下著名论断："计算机时代随处可见，唯独在生产率的统计数字中不见踪影。"

\begin{definition}[索洛悖论（Solow paradox）]
索洛悖论指数字技术投资高速扩张与总体生产率统计表现平淡并存的反常现象：技术普及的速度与生产率统计的反映之间出现系统性落差。
\end{definition}

悖论的核心张力在于：信息技术普及的速度与生产率统计的平淡表现形成显著反差。\mn{索洛悖论原文（1987 年 7 月，索洛为《纽约时报》撰写的书评）："You can see the computer age everywhere but in the productivity statistics." 三个解释维度：测度问题、时滞效应、分配效应。论述题按此框架展开。}这一观察开启了延续三十余年的生产率悖论争论，三种解释（8.2.2 节）分别指向统计测度、技术扩散与收益分配三个层面。

\subsection{悖论的三种解释}

理解索洛悖论的关键在于确定缺口发生的环节：统计是否漏记、红利是否迟到、收益是否被少数主体攫取。三种解释分别回答了这三个追问，也对应三套不同的政策含义：改进统计、耐心等待、调整分配。

\textbf{解释一：测度问题。}生产率统计系统性遗漏了数字经济的产出。数字产品的质量改进（更快、更强的设备价格反而下降）使价格指数高估通胀、低估实际产出；免费数字产品（搜索、社交、导航）的消费者剩余完全不计入 GDP；无形资本（数据、软件、组织资本）的统计核算滞后。三个遗漏的方向一致，\important{都使数字经济时代的产出增长被低估}。其中质量调整遗漏的方向性可以通过一个简单模型明确：统计记录的是名义价格，消费者实际购买的是质量加权的服务价格，两者之差就是漏记的福利增长。

\begin{derivation}{质量调整的方向性}{}
\textbf{质量调整的方向性论证。}设产品名义价格不变（$p_t = p_0$），但产品质量提升：单位产品提供的服务流量从 $q_0$ 增至 $q_t = q_0(1+g)^t$。消费者真正购买的是服务流量，单位流量的"服务价格"为名义价格除以质量 $p_t/q_t$。对该价格取对数求导，分解名义价格变化与质量变化：
\[
\frac{d\ln(p_t/q_t)}{dt} = \frac{d\ln p_t}{dt} - \frac{d\ln q_t}{dt} = 0 - \ln(1+g) \approx -g,
\]
即真实服务价格每年下降 $g$。价格指数若未做质量调整，记录的只是名义价格变化（通胀为零），质量改进带来的真实价格下降被完全漏记。
\end{derivation}

消费者以不变支出获得更多服务：真实服务价格每年下降 $g$ 而统计通胀为零，真实福利增长 $g$ 未反映在统计中。图~\ref{fig:qualityadj} 给出了这一偏差的直观示意：统计价格指数是一条水平线（记录通胀为零），真实服务价格逐年下滑，两条线张开的阴影区就是被统计低估的真实福利增长。对于数字产品，$g$ 由摩尔定律（第 1 章）刻画，计算能力每 18 个月翻番，质量调整的遗漏量随技术迭代累积，统计偏差随时间指数放大。

\begin{figure}[htbp]
\centering
\begin{tikzpicture}[>=Stealth, scale=1]
\fill[main!8] (0,0) -- (6,0) -- (6,-1.2) -- cycle;
\draw[->] (0,-2.6) -- (0,0.9) node[above left] {$\ln$ 价格指数};
\draw[->] (0,0) -- (6.6,0) node[below left=8pt] {$t$};
\draw[main,thick] (0,0) -- (6,0);
\draw[second,thick,dashed] (0,0) -- (6,-1.2) node[pos=0.45, sloped, below, font=\small] {真实服务价格年降 $g$};
\node[anchor=west,font=\small] at (2.0,0.35) {统计价格指数（通胀为零）};
\end{tikzpicture}
\caption{质量调整遗漏的方向}
\label{fig:qualityadj}
\end{figure}

测度问题的第二个来源是免费数字产品。搜索、社交、导航等服务的货币价格为零，其消费者剩余完全不计入 GDP；讨论这一遗漏的规模，需要先厘清消费者剩余的两种度量口径。

\begin{definition}[马歇尔剩余与希克斯剩余（Marshallian and Hicksian surplus）]
马歇尔消费者剩余（Marshallian consumer surplus）以普通需求曲线为基准：价格为 $p$ 时，需求曲线下方、价格线上方的面积度量消费者支付意愿与实际支出之差。希克斯消费者剩余（Hicksian consumer surplus）以补偿需求曲线为基准：收入随价格调整以保持效用水平不变，对应需求曲线下方的面积；按效用基准不同分为补偿变化（CV，以新效用水平为基准）与等价变化（EV，以原效用水平为基准）。两者差异来自收入效应：价格变化改变实际收入，普通需求曲线混同收入效应与替代效应，补偿需求曲线只保留替代效应。
\end{definition}

免费数字产品的消费者剩余完全由支付意愿构成，规模巨大而统计记录为零，这一遗漏的方向同样是产出与 TFP 低估：消费者福利在增长，而 $Y$ 与余值 $\dot A/A$ 的统计记录原地不动。GDP-B 是 Brynjolfsson 团队为此提出的核算尝试：用大规模在线选择实验直接估计用户放弃某项免费服务所要求的补偿（放弃使用的意愿接受补偿，willingness to accept）。代表性结果为：美国中位数用户每月需约 48 美元的补偿才愿放弃 Facebook，按此口径免费数字产品的福利贡献使美国经济增长率每年上调约 0.05--0.11 个百分点；智能手机相机质量改进的福利贡献约为每年 0.63 个百分点（Brynjolfsson 团队，2025）。\mn{hedonic 回归（特征价格法）简释：把产品价格对其性能特征（处理器速度、存储容量等）做回归，回归系数即各特征的隐含价格，统计机构据此分离价格变化中的质量改进成分，如美国劳工统计局自 2018 年起对智能手机采用 hedonic 调整。}该团队研究的结论链构成测度问题解释的实证脉络：企业层面的早期研究已发现 IT 投资回报显著，总量统计却"看不见"，引入质量调整与免费产品福利后缺口大幅收窄；测度问题由此成为悖论解释中实证基础最扎实的部分。

\textbf{解释二：时滞效应。}数字技术的生产率红利存在显著滞后，最有影响的解释来自通用目的技术假说。

\begin{definition}[通用目的技术（GPT）]
通用目的技术（general purpose technology）指可广泛应用于众多行业、持续引发互补性创新并重塑生产方式的技术，蒸汽机、电力、互联网与人工智能皆属此类；其生产率红利的显现依赖配套的互补性投资（组织变革、技能更新），因而存在以十年计的红利时滞。
\end{definition}

电力是历史参照：电动机发明数十年后，工厂重组（从中央动力轴到分布式电机）完成，生产率红利才全面显现。数字技术同样要求互补性投资——组织流程再造、员工技能更新、商业模式重构——互补投资完成之前，技术本身不转化为生产率。这一机制可以形式化。

\begin{derivation}{互补投资与时滞}{}
设新技术 $T$ 的生产率效应通过互补资本 $K_c$（组织资本、技能资本）实现，生产函数为
\[
Y = A_0 \cdot T^{\alpha} \cdot K_c^{\gamma}, \qquad \alpha > 0,\; \gamma > 0.
\]
互补资本按储蓄率 $s_c$ 从产出中积累：
\[
\dot K_c = s_c Y - \delta K_c.
\]
技术 $T$ 在一次投资后即刻到位，但 $K_c$ 只能逐步积累，技术的边际产出
\[
MP_T = \frac{\partial Y}{\partial T} = \alpha A_0 T^{\alpha-1} K_c^{\gamma}
\]
在 $K_c$ 积累完成之前处于低位，只有当 $K_c$ 达到与新技术匹配的水平后，$MP_T$ 才显现。
\end{derivation}

技术到位与生产率红利显现之间的时间差，即 GPT 时滞，其长度取决于互补资本的积累速度（$s_c$）与折旧（$\delta$）：组织变革越慢（$s_c$ 小）、旧组织惯例越顽固（$\delta$ 大），时滞越长。\mn{GPT 假说的政策含义（论述题素材）：数字基础设施投资不能指望立竿见影的生产率回报；企业数字化转型的成效取决于组织变革与技术部署的匹配（第 11 章数字化转型决策与此呼应）。}这一模型同时给出索洛悖论的政策推论：统计上"投资激增而 TFP 平淡"的悖论区间，恰是互补资本积累的过渡期；判断数字技术是否"真的无效"，需要等到 $K_c$ 的收敛。以季度为单位的生产率统计，天然低估了以十年为周期的 GPT 红利。图~\ref{fig:gptlag} 给出了这一时滞的直观示意：技术投资在 $t_0$ 一步到位（阶跃实线），TFP 却随互补资本积累缓慢爬升（虚线），$t_0$ 到 $t_1$ 的阴影区即"投资激增而 TFP 平淡"的悖论窗口，其中 $t_1$ 取 TFP 接近新水平九成以上的时点。

\begin{figure}[htbp]
\centering
\begin{tikzpicture}[>=Stealth, scale=1]
\fill[main!8] (1.2,0) rectangle (4.9,5.2);
\draw[->] (0,0) -- (7.2,0) node[below left=8pt] {$t$};
\draw[->] (0,0) -- (0,5.4);
\draw[main,thick] (0,1) -- (1.2,1);
\draw[main,thick] (1.2,4.2) -- (7,4.2);
\draw[main,dotted,thick] (1.2,1) -- (1.2,4.2);
\draw[second,thick,dashed] (0,1) -- (1.2,1);
\draw[second,thick,dashed] plot[domain=1.2:7,samples=80] (\x,{1+3.2*(1-exp(-0.8*(\x-1.2)))});
\draw[dashed,structurecolor] (1.2,0) -- (1.2,5.2);
\draw[dashed,structurecolor] (4.9,0) -- (4.9,5.2);
\node[below,font=\small] at (1.2,0) {$t_0$};
\node[below,font=\small] at (4.9,0) {$t_1$};
\node[font=\small] at (3.05,1.3) {悖论区间};
\draw[main,thick] (7.4,5.0) -- (8.0,5.0);
\node[anchor=west,font=\small] at (8.1,5.0) {技术投资 $T$};
\draw[second,thick,dashed] (7.4,4.6) -- (8.0,4.6);
\node[anchor=west,font=\small] at (8.1,4.6) {全要素生产率 TFP};
\node[above=2pt,font=\small] at (1.2,5.2) {技术投资到位};
\node[above=2pt,font=\small] at (4.9,5.2) {TFP 充分显现};
\node[font=\footnotesize\color{structurecolor}] at (3.05,0.35) {互补资本积累期};
\end{tikzpicture}
\caption{GPT 时滞}
\label{fig:gptlag}
\end{figure}

\textbf{解释三：分配效应。}数字技术的收益在主体间高度不均等分布：赢家通吃结构（第 3、6 章）使生产率的收益集中于少数超级明星企业，而统计上体现为"平均生产率"的增长乏力；同时，若数字技术的收益以消费者剩余形式存在（免费产品），则国民账户完全无法捕捉。1995 年之后美国生产率回升、并在 2000 年代后再度放缓的实证轨迹（案例~\ref{case:prodrevival}），被解释为测度问题与时滞效应共同作用的结果，索洛悖论至今并无单一裁决，三种解释相互补充。

\begin{case}[美国 1995 年后的生产率回升]\label{case:prodrevival}
索洛悖论提出后最有力的反证来自美国 1990 年代后半期的数据。1973–1995 年美国非农部门劳动生产率年均增速仅约 1.4\%，是"生产率放缓"争论的直接背景；1995–2000 年这一增速回升至年均接近 3\%，为 1970 年代以来最快的五年。对美国数据的增长核算显示，这次回升的主要来源是 IT 资本深化与半导体行业的技术进步：计算机价格暴跌使企业大规模更新 IT 设备，而消化这些设备所必需的互补性组织调整在 1990 年代中后期基本完成。这一轨迹被解读为对 GPT 时滞假说的直接支持：数字技术的生产率红利并未缺席，只是迟到了约二十年；2000 年代后增速再度放缓，则与测度问题及其他结构性因素有关。
\end{case}

\begin{misconception}
\textbf{误区}："索洛悖论意味着数字技术对增长无效"。悖论是统计与理论的紧张关系，而非技术无效的证据：三种解释分别指向测度遗漏、红利时滞与收益分配，数字技术的增长贡献可能真实存在但未被统计捕捉（测度问题）、尚未兑现（时滞效应）或被少数主体攫取（分配效应）。把悖论直接等同于"技术无效"属于常见误读。
\end{misconception}

\subsection{鲍莫尔病与数字化的结构性出路}

索洛悖论讨论的是数字技术自身为何"看不见"，另一个结构性问题从相反方向提问：当经济体日益服务化时，服务业整体的低生产率是否会拖累长期增长？这就是鲍莫尔病，数字技术被认为提供了潜在的出路，其机制构成一道常考名词解释与简答。

\begin{definition}[鲍莫尔病（Baumol's cost disease）]
鲍莫尔病是指经济体中生产率停滞部门的工资随生产率进步部门上升，导致停滞部门的相对成本与价格持续上涨、产出份额扩张，从而拖累整体生产率增速的现象。
\end{definition}

机制上，经济分为生产率快速进步的部门（制造业）与停滞部门（教育、医疗、表演艺术等面对面服务）。劳动力市场一体化使两类部门的工资同涨，而停滞部门的产出无法用资本替代劳动，单位成本与相对价格持续上升；消费者对停滞部门产品的需求价格弹性低，其就业与名义支出份额反而扩张，整体生产率增速被稀释。

数字技术对鲍莫尔病的作用通过三条渠道实现。其一，\textbf{结构红利}：数字产品服务业（软件、云服务）本身是高生产率的生产性服务业，其扩张直接改善服务业的部门构成；其二，\textbf{就业替代}：AI 与自动化替代低生产率服务岗位（客服、初级文书），降低停滞部门的劳动投入权重；其三，\textbf{数字化改造}：零边际成本与规模经济打破面对面服务的产出约束，在线教育、远程医疗使停滞部门首次获得规模经济。同时存在反方向的力量：数字部门与传统服务部门的生产率差距扩大，劳动力若无法跨部门流动，鲍莫尔病反而加剧，数字化的净效应取决于劳动力再配置的速度。\mn{鲍莫尔病与数字经济的关系题：先答定义（工资同涨、相对价格上升、产出份额扩张），再答三条渠道（结构红利、就业替代、数字化改造），最后加一句"净效应取决于劳动力再配置速度"。}对数字经济增长的意义在于：\important{数字化的增长贡献不仅体现在数字部门自身，更体现在对停滞服务部门的改造上}，这是产业数字化（第 11 章）的宏观价值所在。

\section{数据要素与增长模型}

\subsection{三要素增长核算}

8.2 节的三种解释都试图在既有增长框架内部消化数字技术的增长之谜；本节转向一个更根本的问题：当数据作为独立的生产要素（而非改进效率的工具）进入生产函数时，增长本身的结构会怎样变化。第 7 章的四属性对照已经说明数据与资本、劳动的本质差异：非竞争性使共享不减少单主体价值，部分排他性使产权执行成本高，复制边际成本趋零使供给激励特殊，价值依赖场景使定价困难。这些微观属性在宏观层面意味着，数据对增长的作用方式不能沿用传统要素的分析框架：资本的使用具有排他性（一台机器同时只能用于一处），劳动同样不可分身，唯有数据可被无限主体同时使用；同时数据的价值又随时间迅速贬值。把 $D$ 引入生产函数，就是把第 7 章的微观制度分析翻译成宏观增长语言。本节先给出静态的三要素核算，8.3.2 与 8.3.3 节再进入增长理论。

\begin{derivation}{三要素增长核算}{}
\textbf{三要素生产函数与核算。}扩展生产函数为
\[
Y = A K^\alpha L^\beta D^\gamma,
\]
其中 $D$ 为数据要素投入。对数微分得
\[
\frac{\dot Y}{Y} = \frac{\dot A}{A} + \alpha\frac{\dot K}{K} + \beta\frac{\dot L}{L} + \gamma\frac{\dot D}{D}.
\]
\end{derivation}

三要素核算把数据纳入增长分解，同时暴露两个新问题。其一，数据贡献 $\gamma \dot D/D$ 的度量难点：需要度量数据存量的增长，而数据的"折旧"对应\textbf{时效衰减}（过时数据的价值衰减）而非物理损耗，其速率由技术与市场结构内生决定，这是数据要素在宏观层面独有的度量障碍。三要素核算是 Jorgenson 增长核算框架向数据要素的直接推广；数据存量的统计口径（数据条数、字节量还是质量加权）尚无统一标准，国家统计局《数字经济及其核心产业统计分类（2021）》给出了行业统计框架，宏观数据存量的测度仍在探索之中。

其二，若 $\alpha+\beta+\gamma>1$，数据的引入使生产呈现规模报酬递增（increasing returns to scale）。所谓规模报酬，指所有要素同时同比例扩大时产出放大的倍数：把 $K$、$L$、$D$ 都扩大 $\lambda$ 倍，产出变为 $\lambda^{\alpha+\beta+\gamma}Y$。若 $\alpha+\beta+\gamma=1$，要素翻番产出正好翻番，这是索洛模型的标准情形；若 $\alpha+\beta+\gamma>1$，要素翻番产出的增加超过一倍，每多投入一份要素，回报不降反升。

规模报酬递增看似是"好事"，却摧毁了竞争性均衡的存在基础。竞争性均衡要求企业按要素的边际产品支付报酬，并且产出恰好能够分配给全部要素；8.1.1 节的欧拉定理保证这个账目只在规模报酬不变时结清（此时各要素报酬之和恰好等于 $Y$）。规模报酬递增下，各要素按边际产品取得的报酬之和为 $(\alpha+\beta+\gamma)Y>Y$，产出不够分，账目无法结清。换个角度，给定要素价格，企业把生产规模扩大 $\lambda$ 倍，成本增长 $\lambda$ 倍、收益却增长 $\lambda^{\alpha+\beta+\gamma}>\lambda$ 倍，利润随规模无休止上升，最优产量不存在。均衡无从谈起，增长的动力学也与基准情形有本质差异：规模报酬不变下增长收敛于索洛稳态、长期增长依赖外生的 $g$；规模报酬递增下积累的收益随规模自我放大，增长可能自我维持、无须外生推力。8.3.2 节的 AK 模型正是对这一差异的形式化。

\subsection{数据要素与内生增长}

数据的非竞争性（第 7 章）在内生增长（endogenous growth）框架中具有革命性含义：传统内生增长理论（Romer 的知识积累、Lucas 的人力资本）中，规模报酬递增的来源同样是"非竞争性投入"（知识、思想），数据不过是这一逻辑的最新载体。

\begin{definition}[AK 模型]
AK 模型是内生增长理论最简的形式化框架：生产函数 $Y = AK$（$A$ 为常数），资本（广义，包含知识）的边际产出恒为 $A$ 而不递减，增长由资本积累内生驱动、无须依赖外生技术进步；"AK"的名称即来自这一线性生产函数形式。
\end{definition}

AK 模型的全部结论可以从积累方程直接写出，关键是观察增长率如何与资本水平脱钩，这是它与索洛模型的本质区别。

\begin{derivation}{AK 模型与数据}{}
\textbf{AK 模型与数据。}设 $Y = AK$（$A$ 为常数），资本积累方程为
\[
\dot K = sAK - \delta K.
\]
两边同除以 $K$，得到资本存量增长率
\[
\frac{\dot K}{K} = sA - \delta,
\]
它与资本水平无关，恒为常数 $g = sA - \delta$。若劳动人口以速率 $n$ 增长，人均资本 $k = K/L$ 的增长率还需扣除人口稀释，人均量的增长率为
\[
g = sA - (n+\delta).
\]
\end{derivation}

将数据视为广义资本的一部分（$K$ 中纳入数据存量 $D$），数据的非竞争性使生产函数的线性假设获得微观基础：数据可被无限主体同时使用，其边际产出不因"使用"而递减，这与实体资本的物理属性形成鲜明对照。AK 模型的增长机制是一个正反馈循环：今天多积累一点资本，明天产出就更高，从而支持更高的储蓄与再投资，只要 $sA>\delta$，循环就不会衰减。实体资本做不到这一点，因为每多一台机器对产出的贡献递减，循环终会停下来；\mn{AK 对比索洛（辨析常客）：$g = sA - \delta$ 中储蓄率 $s$ 直接进入长期增长率；索洛稳态里 $s$ 只改变水平，长期增长率只看 $g$。}数据、知识这类非竞争性投入则不同，边际产出不递减使循环可以自我维持。这一理论含义是：\important{数据要素具有将增长从"外生技术进步驱动"推向"内生积累驱动"的潜力}；但这一潜力的兑现取决于制度条件（数据的流动与共享），数据孤岛使数据的非竞争性无法实现，增长潜力被产权与合规成本锁死，构成第 7 章堵点的宏观代价。

\textbf{Jones–Tonetti 数据所有权模型}（Jones and Tonetti, 2020, AER, "Nonrivalry and the Economics of Data"）进一步考察了数据产权配置的增长含义，其核心权衡可以用一个简化的形式化框架阐明。

\begin{derivation}{数据产权配置}{}
\textbf{数据产权配置的增长比较。}设经济中有 $n$ 个同质企业、一个消费者（数据主体）。消费者的数据存量 $d$ 若被企业 $i$ 使用，该企业的生产率提升为 $A(d)$（$A' > 0$，$A'' < 0$，且 $A(0) = 0$：没有数据时生产率为零）。数据具有非竞争性：多个企业可同时使用同一数据而不互相妨碍。企业的产出为 $y_i = A(d_i)\, k_i^\alpha$，其中 $d_i$ 为企业实际使用的数据量。数据由消费者生成，生成成本为隐私成本 $e(d)$（$e' > 0$）：无论数据最终由谁持有，这一成本都要付出，区别仅在于由谁承担、采集激励如何传导。

\textbf{情形一：消费者拥有数据。}消费者向每个企业出售数据使用权，定价 $p$。每个企业使用全部数据（$d_i = d$），总产出为
\[
Y_C = \sum_{i=1}^{n} A(d)\, k_i^\alpha = n \cdot A(d)\, k^\alpha \qquad (\text{对称情形}).
\]
消费者获得的许可收入为 $n p$，采集激励取决于 $n p$ 与隐私成本 $e(d)$ 的权衡：最优数据供给满足
\[
n p = e'(d).
\]
若许可市场完善（数据可定价、交易成本低），每个企业的许可价格等于数据对其自身的边际产出 $A'(d) k^\alpha$，许可总收入 $np = nA'(d) k^\alpha$，上式恰为社会最优条件。

\textbf{情形二：企业拥有数据。}单一企业独占数据（$d_1 = d$，$d_i = 0,\; i \geq 2$），总产出为
\[
Y_F = A(d)\, k^\alpha < Y_C.
\]
非竞争性资源被排他使用，社会产出损失为 $(n-1)A(d)k^\alpha$，这是数据独占的静态效率损失。但企业拥有的采集激励更强：独占者将数据对其自身使用的价值内部化，为获得数据承担同样的生成成本 $e(d)$，最优采集满足
\[
A'(d) k^\alpha = e'(d).
\]
对照社会最优条件 $nA'(d)k^\alpha = e'(d)$，独占者只内部化社会边际价值的 $1/n$，采集相对社会最优仍不足。
\end{derivation}

两种产权配置的比较因此是明确的：消费者拥有实现了非竞争性的效率（数据被充分共享），但依赖数据交易市场传递采集激励；企业拥有保障了采集激励，但以共享损失为代价。交易市场的完善程度决定消费者拥有方案的实际表现：若数据难以定价、交易成本高（第 7 章三大堵点中的"定价难"），许可价格无法形成或严重偏低，消费者失去采集激励，数据供给不足，这正是第 7 章 Samuelson 条件的结论在产权配置上的再现：私人供给激励不足使公共品性质的数据供给不足；相对许可市场失灵的情形，企业拥有的激励显著更强，但共享损失依然存在。\important{最优产权设计的目标是分离两种激励}：使"采集"与"共享"在不同权利层级上分别获得回报。中国数据二十条的三权分置（第 7 章）正是这一分离的制度实践：持有权与加工使用权的配置照顾采集激励，产品经营权的流通实现共享收益。

\subsection{数据积累与稳态}

8.3.1 的三要素核算是静态分解，8.3.2 说明数据改变了增长的发动机；两者都未回答数据存量本身如何随时间演化。将第 7 章的数据时效性（价值随时间衰减）纳入积累方程，可以考察数据要素的长期动态。

\begin{derivation}{数据积累与稳态}{}
\textbf{数据积累与稳态增长。}设生产函数 $Y = A K^\alpha L^\beta D^\gamma$，其中 $\gamma < 1$（数据对产出的规模弹性小于 1；若 $\gamma \geq 1$，稳态不存在，产出将随数据积累而爆炸式增长）。数据存量的积累方程为
\[
\dot D = s_D Y - \delta_D D,
\]
其中 $s_D$ 为"数据储蓄率"（产出中投入数据采集与治理的份额），$\delta_D$ 为数据折旧率（时效衰减）。稳态要求 $\dot D = 0$，得
\[
D^* = \frac{s_D}{\delta_D} Y^*.
\]
将其代入生产函数消去 $D$（即在稳态条件 $\dot D = 0$ 下解出稳态产出 $Y^*$）：
\[
Y^* = A (K^*)^{\alpha} (L^*)^{\beta} \left(\frac{s_D}{\delta_D} Y^*\right)^{\gamma}
\;\Rightarrow\;
Y^* = A^{\frac{1}{1-\gamma}} (K^*)^{\frac{\alpha}{1-\gamma}} (L^*)^{\frac{\beta}{1-\gamma}} \left(\frac{s_D}{\delta_D}\right)^{\frac{\gamma}{1-\gamma}}.
\]
\end{derivation}

稳态的弹性放大同样可以借助"池子"来理解：采集投资 $s_D Y$ 是入水，时效衰减 $\delta_D D$ 是漏水，稳态是两者平衡处；稳态产出对技术、资本与劳动的弹性被放大为原来的 \important{$1/(1-\gamma)$} 倍：$\gamma$ 接近 1 时（数据对产出的贡献趋强），稳态产出对技术、资本与劳动的弹性急剧放大，\important{数据要素使规模报酬递增成为稳态的常态而非例外}。\mn{数据稳态速记：稳态条件 $\dot D=0$ 得 $D^*=\frac{s_D}{\delta_D}Y^*$；回代生产函数得弹性放大 $\frac{1}{1-\gamma}$，$\gamma$ 越接近 1 放大越猛。}同时，数据折旧率 $\delta_D$ 以幂次 $\gamma/(1-\gamma)$ 进入稳态：时效衰减越快，数据要素的长期贡献越受侵蚀，数据治理（保鲜、更新、维护）因此具有不亚于数据采集的增长含义。

\begin{misconception}
\textbf{误区}："数据越多，增长越快"。积累方程与稳态显示，数据对增长的贡献同时受采集率 $s_D$ 与折旧率 $\delta_D$ 的约束：时效衰减使部分采集投入不断被抵销，数据孤岛（第 7 章）使已采集的数据无法进入生产函数，数据越多未必增长越快；数据治理与流通制度的效率才是把 $D$ 转化为增长的关键。
\end{misconception}

\section{数字基础设施与长期增长}

\subsection{基础设施的外部性}

8.3 节的模型把数据视为企业可自由积累与使用的要素，但数据的采集、流动与计算依赖网络、算力等基础设施，基础设施供给不足时，数据要素的增长潜力无从兑现。数字基础设施（digital infrastructure，网络、算力、数据中心）具有公共品与正外部性双重属性：单个企业的投资决策不考虑其对其他主体生产率的外溢收益，私人投资低于\important{社会最优}。

\begin{derivation}{基础设施外部性}{}
\textbf{外部性与增长含义。}设数字基础设施存量 $G$ 进入各企业的生产函数：
\[
y_i = A(G)\, k_i^\alpha,
\]
其中 $A'(G) > 0$。单个企业投资基础设施的收益为 $\partial y_i/\partial G$（私人边际收益），而基础设施的社会边际收益为
\[
\sum_i \frac{\partial y_i}{\partial G} = \sum_i A'(G) k_i^\alpha,
\]
即全体使用者的边际收益之和，私人收益与社会收益的缺口为
\[
\sum_{i \neq j} A'(G) k_i^\alpha > 0.
\]
\end{derivation}

这一缺口是正外部性的度量：\important{私人投资激励不足，公共投资的效率依据}。\mn{论述题模板：数字基础设施为何要公共投资？私人收益 $A'(G)k_i^\alpha$ 小于社会收益 $\sum_i A'(G)k_i^\alpha$，正外部性缺口致私人投资不足，故公共投资有据；东数西算、5G、千兆光网皆为例证。}这为"东数西算"工程、5G 网络、千兆光网等国家数字基础设施战略提供了标准的经济学论证，其逻辑与第 7 章 Samuelson 条件同源：非竞争性投入的边际社会价值是各主体边际价值之和。

\begin{case}[东数西算]
"东数西算"（channeling computing resources from the east to the west）工程是中国算力基础设施空间布局的国家战略。2021 年 12 月《"十四五"数字经济发展规划》提出实施"东数西算"工程；2022 年 2 月，国家发展改革委等部门批复同意在京津冀、长三角、粤港澳大湾区、成渝、内蒙古、贵州、甘肃、宁夏 8 地启动国家算力枢纽节点建设，规划 10 个国家数据中心集群。其经济逻辑在于：东部产生的数据量占全国大半，但土地、电价与能耗指标紧张；西部绿色能源富集、气候凉爽，适宜数据中心布局。将东部产生的数据引导至西部计算，既降低了全社会算力供给的成本，也通过算力枢纽向西部布局带动当地数字产业发展，缩小区域数字鸿沟（8.4.2 节）。在本章框架下，东数西算正是基础设施外部性的公共投资回应：私人投资不足的算力设施，由国家统筹补齐。
\end{case}

\subsection{数字鸿沟与增长收敛}

数字鸿沟（digital divide）指不同群体与地区在数字接入与使用上的系统性不平等（第 9 章详述）。在增长理论中，数字鸿沟作用于收敛过程：接入差距改变各经济体的技术吸收能力与稳态位置。\mn{$\beta$ 收敛口诀：$\beta$ 收敛看速度（负相关），条件收敛看控制（控住稳态才成立），绝对收敛不成立。}本节先辨析收敛的概念，再推导收敛方程，最后讨论数字技术对收敛的双重作用。

\begin{definition}[$\beta$ 收敛（beta convergence）与条件收敛（conditional convergence）]
$\beta$ 收敛指人均产出增长率与初始水平负相关：初始位置离自身稳态越远，增速越高。条件收敛是 $\beta$ 收敛的实证版本：只有控制住储蓄率、人口增长率、技术水平等稳态决定因素（让比较对象拥有相同的稳态 $y^*$）后，落后经济的增速才更高；不控制稳态差异的绝对收敛（unconditional convergence）在跨国数据中不成立。数字鸿沟之所以进入增长分析，正是因为它改变的就是各经济体的稳态位置。
\end{definition}

\begin{derivation}{β 收敛}{}
\textbf{$\beta$ 收敛方程。}由 8.1.2 节的动态方程 $\dot k = s k^\alpha - (\delta+n+g)k$ 出发，两边同除以 $k$ 得
\[
\frac{\dot k}{k} = s k^{\alpha-1} - (\delta+n+g).
\]
利用 $d\ln y/dt = \alpha\, d\ln k/dt$，令 $x = \ln k$，在 $x^* = \ln k^*$ 附近作一阶泰勒展开：
\[
\frac{d\ln y}{dt} = \alpha\Bigl[s e^{(\alpha-1)x} - (\delta+n+g)\Bigr] \approx \alpha s (\alpha-1) e^{(\alpha-1)x^*}\,(x - x^*).
\]
利用 $e^{(\alpha-1)x^*} = (k^*)^{\alpha-1}$ 与稳态条件 $s (k^*)^{\alpha-1} = \delta+n+g$，上式系数化简为 $-(1-\alpha)\,\alpha\,(\delta+n+g)$；再由 $\ln y - \ln y^* = \alpha(x - x^*)$，得
\[
\frac{d\ln y}{dt} \approx -\beta\bigl(\ln y - \ln y^*\bigr), \qquad \beta = (1-\alpha)(\delta + n + g).
\]
落后经济体（$\ln y < \ln y^*$）增速更高，向稳态收敛，收敛速度由 $\beta$ 度量。
\end{derivation}

图~\ref{fig:betaconv} 给出了这一双重作用的示意：两条收敛路径中，落后经济的初始增速更高（曲线更陡，这是 $\beta$ 收敛），但数字鸿沟压低了它的稳态目标线，收敛完成后收入差距依然存在。数字技术对收敛的双重作用。一方面，数字技术是通用技术，落后地区可以跨越式采用（移动支付跳过信用卡阶段），直接吸收最新技术而不必重复先行者的全部技术阶梯，收敛加速；另一方面，数字鸿沟使落后地区的有效技术 $A$ 系统性偏低：若数字接入是生产函数的一部分（$A$ 依赖于 $G$），则基础设施差距压低的是落后地区的\textbf{稳态} $y^*$，给定当前的 $\ln y$，向稳态收敛的距离 $\ln y - \ln y^*$ 被拉大，落后地区将向一个更低的稳态收敛，稳态收入差距被持久锁定；若接入差距还拖累新技术的吸收（技术进步率 $g$ 更低），则收敛速度 $\beta$ 本身下降，分化进一步固化。实证研究（世界银行《2016 年世界发展报告：数字红利》）的共识是：\important{数字技术的收益是条件性的}：只有当互补条件（人力资本、制度质量、基础设施）满足时，数字技术才促进收敛；否则可能加剧分化。

\begin{figure}[htbp]
\centering
\begin{tikzpicture}[>=Stealth, scale=1]
\draw[->] (0,0) -- (9.2,0) node[below left=8pt] {$t$};
\draw[->] (0,0) -- (0,7.2) node[above left] {$\ln y$};
\draw[dashed,structurecolor] (0,5.2) -- (8.0,5.2);
\draw[dashed,structurecolor] (0,3.4) -- (8.0,3.4);
\node[anchor=east,font=\small] at (0,5.2) {$\ln y^*_H$};
\node[anchor=east,font=\small] at (0,3.4) {$\ln y^*_L$};
\draw[main,thick] (0,5.8) -- (1,5.564) -- (2,5.421) -- (3,5.32) -- (4,5.262) -- (5,5.228) -- (6,5.212) -- (7,5.204) -- (8,5.199) node[pos=0.3, sloped, above=6pt, font=\small] {发达地区（缓）};
\draw[second,thick] (0,2.2) -- (1,2.672) -- (2,2.958) -- (3,3.151) -- (4,3.272) -- (5,3.347) -- (6,3.392) -- (7,3.417) -- (8,3.431) node[pos=0.3, sloped, below=6pt, font=\small] {落后地区（初速更高，$\beta$ 收敛）};
\draw[->,thick] (6.5,5.12) -- (6.5,3.48);
\node[anchor=west,font=\small,align=left] at (6.62,4.3) {数字鸿沟\\压低稳态};
\draw[<->,second,thick] (8.05,5.2) -- (8.05,3.4) node[midway,right=4pt,font=\small,align=left]{稳态\\差距\\锁定};
\end{tikzpicture}
\caption{$\beta$ 收敛与数字鸿沟}
\label{fig:betaconv}
\end{figure}

\begin{chapterbox}
\textbf{本章考点清单}

\begin{enumerate}
\item 增长核算方程的对数微分推导、欧拉定理与索洛余值 \important{【学硕·热点】【专硕·核心】}
\item 要素贡献率分解的计算方法（数值算例）\important{【学硕·热点】}
\item 索洛稳态 $k^* = (s/(\delta+n+g))^{1/(1-\alpha)}$ 与技术进步决定长期增长 \important{【学硕·热点】}
\item 黄金律水平 $k_{GR} = (\alpha/(\delta+n+g))^{1/(1-\alpha)}$ 与动态无效率判断 \important{【学硕·热点】}
\item 索洛悖论及其三种解释（测度、时滞、分配）\important{【专硕·核心】}
\item 质量调整遗漏的方向性证明 \important{【专硕·核心】}
\item 消费者剩余测度之辨（马歇尔 vs 希克斯）与 GDP-B \important{【专硕·扩展】}
\item 鲍莫尔病与数字化的三条出路（结构红利、就业替代、数字化改造）\important{【专硕·扩展】}
\item 三要素生产函数 $Y = AK^\alpha L^\beta D^\gamma$ 的增长核算 \important{【专硕·核心】}
\item 数据非竞争性与 AK 模型的内生增长含义 \important{【专硕·核心】}
\item Jones–Tonetti 所有权模型的核心权衡 \important{【专硕·扩展】}
\item 数据积累稳态、弹性放大 $1/(1-\gamma)$ 与数据存量测度争议 \important{【专硕·扩展】}
\item 数字基础设施的外部性论证（私人 vs 社会收益）\important{【专硕·核心】}
\item $\beta$ 收敛方程与数字鸿沟对收敛的双重作用 \important{【专硕·扩展】}
\end{enumerate}
\end{chapterbox}

增长核算与增长理论回答的是"蛋糕能否做大"，但增长的总量效应最终要落到具体的劳动者与家庭身上。数字经济让哪些人的收入更快增长、替代了哪些人的工作、又把哪些群体挡在红利之外——分配维度的问题，是增长分析的必然延续。
